% Audited research construction only; unregistered; no ordinary-Chow survival assertion.
\paragraph{A direct exceptional SNC model with nonzero triple-point periods
(research construction).}
For every prime $p$ we construct an algebraizable normal isolated
threefold germ over an algebraically closed field of characteristic $p$,
a projective regular resolution with three rational SNC exceptional
components and a positive negative-ample exceptional divisor, for which
the period condition of C.3bq is false. This is a failure of that
sufficient condition, not a nonvanishing assertion for the ordinary
Chow quotient.

\emph{1. Explicit constant quadrics and generic remaining coefficients.}
Put $k_0=\overline{\mathbf F}_p$ and use coordinates
$[X_0:X_1:X_2:X_3]$ on $\mathbf P^3$.
If $p\ne2$, set
\[
 A=X_1^2-X_0X_3,\qquad
 B=X_2^2-X_1X_3+X_0X_1,
 \qquad Q_1=B+\iota A,\quad Q_2=B-\iota A,
 \quad \iota^2=-1.
\]
If $p=2$, set
\[
 A=X_1^2-X_0X_3,\qquad
 B=X_2^2+X_0X_2-X_1X_3,
 \qquad Q_1=B,\quad Q_2=A+B.
                                                        \tag{NP.1}
\]
In both cases the quadrics $Q_1,Q_2$ are smooth and their
intersection $C$ is a smooth connected elliptic quartic.
For odd $p$, the determinant of the linear gradient matrix of $B+\lambda A$ is
$-4\lambda(\lambda-1)(\lambda+1)$, so it has no common projective
zero if $\lambda\notin\{0,1,-1\}$; apply this to
$\lambda=\pm\iota$. On $X_0\ne0$, $C=(A=B=0)$ is
$y^2=x^3-x$, while its single point at infinity is
$[0:0:0:1]$, where the two gradients are independent.
For $p=2$, both displayed quadrics have nonsingular gradient matrices,
and the affine elliptic equation is $y^2+y=x^3$; the same infinity
check applies. These verifications use no characteristic-zero
specialization. Smooth connected quadrics over $k_0$ are rational.

Let $t_1,t_2,t_3$ and the $84$ symbols $u_I$, indexed by the
monomials $X^I$ of degree six, be algebraically independent over $k_0$.
Over the algebraic closure $k$ of their rational function field, put
\[
 H=X_0+t_1X_1+t_2X_2+t_3X_3,
 \qquad G=\sum_{|I|=6}u_I X^I.
                                                        \tag{NP.2}
\]
The tangent-cone divisor
\[
 E_0=(Q_1Q_2H=0)\subset\mathbf P^3_k
\]
is SNC. Its three double curves are $C_{12}=C_k$ and the two
smooth conics $C_{13},C_{23}$. Its triple locus consists of four
distinct points $q_1,\ldots,q_4=C\cap H$.
The generic sextic $G$ is nonzero at all four triple points and has
simple zeros on all three double curves. These are simultaneous
nonempty open conditions: apply the first-jet Bertini argument to
$H$ on the fixed smooth quadrics and elliptic curve, and to $G$ on
the finitely many smooth double curves, with the finitely many
nonvanishing conditions at the triple points
\citep[Lemma 33.47.3 and proof]{D_StacksBertiniUniformHigher}.

\emph{2. The affine germ and its first blowup.}
Let
\[
 F=Q_1Q_2H-G\in k[X_0,X_1,X_2,X_3],\qquad
 R_0=\bigl(k[X_0,X_1,X_2,X_3]/(F)\bigr)_{(X_0,X_1,X_2,X_3)}.
                                                        \tag{NP.3}
\]
The degree-five initial form of $F$ is $Q_1Q_2H$, which is squarefree.
Hence the associated graded ring is reduced and so is $R_0$.
Blow up its maximal ideal, writing $Y\to\operatorname{Spec}R_0$.
On each standard blowup chart, with radial coordinate $t$, its
strict-transform equation is
\begin{equation}
                   Q_1(u)Q_2(u)H(u)=tG(u),
             \qquad E_0=(t=0).                        \tag{NP.4}
\end{equation}
The maximal ideal becomes the invertible ideal $(t)$; there are
no extra vertical components in the Rees blowup. This can also be
seen by saturating the chart equation by $t$.

At a smooth point of $E_0$ the equation is regular. At a double
point where $G\ne0$, solve for $t$; the ambient space is regular
and its exceptional divisor has the two transverse branches.
At a triple point, $G$ is a unit; the equation similarly makes
the ambient space regular and the divisor SNC.
At a zero of $G$ on a double curve away from the triple points,
transversality gives étale coordinates in which (NP.4) is exactly
\begin{equation}
                         xy=tz.                        \tag{NP.5}
\end{equation}
Thus the only singularities along the closed fibre are ordinary
threefold nodes: $24$ on $C_{12}$ and $12$ on each conic, for a
total of $48$. The counts are the products of the double-curve
degrees $4,2,2$ with six, and every point is reduced.

Here is a direct local-domain verification, avoiding a geometric
irreducibility assumption about a generic affine polynomial.
Every irreducible component of the blowup meets the closed fibre
by properness over the local base, and the exceptional ideal is
Cartier, so that intersection contains a two-dimensional component
of $E_0$. At a general point of any double curve with $G\ne0$,
the ambient space is regular; hence its two adjacent exceptional
components lie in the same ambient irreducible component.
The triangle of components is connected, so all ambient components
coincide. The blowup, and hence the reduced local ring $R_0$, is
integral. After the resolution in Step 3, properness over the local
base will imply regularity everywhere from regularity along the
closed fibre. Its complement is unchanged from the puncture of
$R_0$. Thus this puncture is regular; the hypersurface $R_0$ is
Cohen--Macaulay and regular in codimension one, hence normal.

\emph{3. Ordered small resolutions, with the marked curves fixed.}
Blow up the Weil component of $E_0$ corresponding to $Q_1$, then
the strict transform of the one corresponding to $Q_2$.
These are projective morphisms. They are isomorphisms where the
chosen component is Cartier, so they affect only the nodes just
listed. At a node whose earlier component is $(x=t=0)$ in
(NP.5), the blowup of $(x,t)$ has saturated charts
\[
 t=xv,\quad y=vz\qquad\text{and}\qquad
 x=tw,\quad z=wy.
                                                        \tag{NP.6}
\]
Both charts are regular in every characteristic. The exceptional
fibre over the node is a single $\mathbf P^1$, contained in the
chosen component; that component is blown up at the point, while
the other component is unchanged. In the first chart the total
exceptional divisor is $xv=0$, so no new divisor or triple stratum
is created. The old double curve is unchanged as an abstract
marked curve. These are also the explicit local charts of
\citet[Lemmas 15--16, published pp. 63--64]{frKollarSNC2014}; the
calculation itself is étale algebra in every characteristic and
uses no discrepancy or analytic-contraction assertion.

Let $q:Z\to Y$ be the composite. Its reduced closed fibre is
$E=E_1\cup E_2\cup E_3$, with
\[
 E_1\simeq\operatorname{Bl}_{36\ \mathrm{points}}(Q_1),
 \qquad E_2\simeq\operatorname{Bl}_{12\ \mathrm{points}}(Q_2),
 \qquad E_3\simeq\mathbf P^2.
                                                        \tag{NP.7}
\]
It is SNC and all its components are rational surfaces.
The double curves and all four marked triple points are those of
$E_0$. All nodes have been resolved. Since $Z$ is proper over the
local base and regular along its entire closed fibre, its closed
nonregular locus has empty proper image and is empty. The same
proper-image observation removes any putative additional
singularity not meeting that fibre. Thus $Z$ is regular and is an
isomorphism over the puncture. Passing to $R=R_0^{\rm sh}$ preserves
all these properties and leaves the exceptional fibre unchanged,
since the residue field is already algebraically closed.

\emph{4. An actual full-support negative ample divisor.}
Write $E=E_1+E_2+E_3$ also for the reduced Cartier divisor.
There are no exceptional divisors of the small resolution, so
$E=q^*E_0$, and
\[
                     \mathcal O_E(-E)=q_E^*\mathcal O_{E_0}(1).
\]
Put $A=2E_1+E_2$. On a small exceptional $\mathbf P^1$ over an
$ij$ node, the line bundle $\mathcal O_Z(-A)$ has degree
$2-1=1$, $2-0=2$, or $1-0=1$, for $ij=12,13,23$ respectively.
Indeed the tautological blowup bundle is $\mathcal O(-E_i)$ for
the earlier component and has degree one; the two component
classes sum to the principal local equation $t$.
Every other fibre of $q$ is a point. Therefore
$\mathcal O_E(-A)$ is $q_E$-ample by fibrewise ampleness for a
proper morphism
\citep[Lemma 37.50.3 and complete proof]{frSNCPeriodFibreAmple}.
Composition of relative ample bundles
\citep[Lemma 29.38.8 and proof]{frC3beStacksCompositeAmple}
now gives, for every sufficiently large integer $M$,
\begin{equation}
 D=M E+A=(M+2)E_1+(M+1)E_2+M E_3,
                   \qquad\mathcal O_E(-D)\ \text{ample}.
                                                        \tag{NP.8}
\end{equation}
Every coefficient is strictly positive. This is the negative
ample combination required in C.3bm and C.3bq. The construction
already maps projectively to the affine hypersurface germ;
neither an analytic contraction nor Artin's contraction theorem
is an input.

\emph{5. The period condition is not automatic.}
The dual complex has three vertices, three edges, and four
triangular faces with the same oriented boundary. Hence
$H_2(\Delta,\mathbf Q)$ has dimension three, generated by
$v_j-v_1$ for $j=2,3,4$. Up to a common orientation sign, their
periods on the elliptic edge $C_{12}$ are
\[
                  [q_j-q_1]\in\operatorname{Pic}^0(C)(k)\otimes\mathbf Q.
\]
They are nonzero. To see this precisely, use a $k_0$-point as
origin on the fixed elliptic curve $C/k_0$.
Every torsion geometric point of this constant elliptic curve
lies in $C(k_0)$, since each torsion subgroup scheme is finite
over the algebraically closed field $k_0$.
The incidence variety of two ordered distinct points of $C$
and a hyperplane through them is a $\mathbf P^1$-bundle over
$(C\times C)\setminus\Delta_C$. It is irreducible, dominates
$C\times C$, and is generically finite over the hyperplane
parameter space. Thus every ordered pair of distinct points in
the geometric generic hyperplane section is generic in
$C\times C$ over $k_0$. It lies on none of the proper closed
graphs $q'=q+c$ with $c\in C(k_0)$. In particular its difference
is nontorsion. Adjoining the independent sextic coefficients
does not change this conclusion. The two conic edges have
zero degree-zero Picard group.

The small-resolution centres avoid all triple points, and the
elliptic double curve is unchanged, so these periods remain
nonzero on the actual exceptional SNC divisor $E$.
This proves that the hypotheses of C.3bq are genuinely additional
even when every exceptional component is rational and a full
negative ample combination exists. It proves nothing about the
ordinary quotient of this particular puncture: component perfect
families and other endpoint-zero corrections might still remove
the corresponding classes. No such survival is asserted.
